\chapter{生产性毒物与职业中毒}
\setcounter{chapter}{3}
\chapterintro{
本章为课程核心章节（12学时）。

\textbf{【教学大纲要求】}

\imp{掌握：}毒物、中毒、生产性毒物和职业中毒的概念；毒物的来源与存在状态；生产性毒物进入人体的途径；毒物在体内的变化过程；影响毒物对机体毒作用的因素\keyword{*}；职业中毒的临床表现、诊断、治疗与预防；铅中毒（理化特性、毒理\keyword{*}、临床表现、诊断、处理原则与预防）；汞中毒（理化特性、毒理\keyword{*}、临床表现、诊断、处理原则与预防）；锰中毒（毒理与临床表现）；刺激性气体（种类、毒作用机制、临床表现——化学性肺水肿\keyword{*}、防治原则）；窒息性气体（概述、分类、毒作用特点\keyword{*}、毒作用表现、治疗原则）；一氧化碳、硫化氢、氰化氢中毒（特性、毒作用机制、临床表现、防治原则）；有机溶剂（理化特性与毒作用特点、对健康的影响）；苯（理化特性、毒理、毒作用机制与临床表现\keyword{*}、诊断、处理原则与预防）；苯的氨基和硝基化合物（毒作用共同点\keyword{*}、临床表现与诊断、处理原则与预防）；高分子化合物（单体与聚合物、远期三致作用）；氯乙烯（特性、毒理、临床表现、诊断与预防）；有机磷酸酯类农药（特性、毒理\keyword{*}、临床表现、诊断、处理原则和预防）。

\imp{熟悉：}常见金属与类金属（铅、汞、锰）；刺激性气体（氯气、氮氧化物、氨、光气）；甲苯、二甲苯、二氯乙烷、二硫化碳的临床表现；苯胺和三硝基甲苯中毒；含氟塑料和丙烯腈中毒；拟除虫菊酯类农药、氨基甲酸酯类农药、百草枯毒作用的特点；三氯乙烯、甲醛中毒。

（注：\keyword{*} 标示教学难点）
}

% ----- 3.1 毒理学基本概念 -----
\section{毒理学基本概念}

\begin{defbox}
\textbf{毒物（Poison）}：在一定条件下，摄入较小剂量即可引起机体暂时或永久性病理改变，甚至危及生命的化学物质。

\textbf{毒素（Toxin）}：生物体（如细菌、真菌、植物、动物等）产生的具有毒性的化学物质。

\textbf{中毒（Poisoning）}：机体受毒物作用后引起一定程度损害而出现的疾病状态。

\textbf{生产性毒物（Productive Toxicant）}：生产过程中产生的，存在于工作环境中的毒物。

\textbf{职业中毒（Occupational Poisoning）}：工作人员在生产过程中由于接触生产性毒物而引起的中毒。
\end{defbox}

\subsection{职业中毒的分类}

\begin{keybox}
职业中毒按照接触毒物种类可分为：
\begin{enumerate}[leftmargin=*]
    \item \imp{金属和类金属}：铅、汞、锰、镉、铍、铊、钒、砷、磷、铀
    \item \imp{刺激性气体}：氯气、光气、氨、氮氧化物
    \item \imp{窒息性气体}：一氧化碳、氰化氢、硫化氢、甲烷
    \item \imp{有机溶剂}：苯、甲苯、二甲苯、二氯乙烷、二硫化碳、正己烷
    \item \imp{苯的氨基和硝基化合物}：苯胺、三硝基甲苯
    \item \imp{农药}：有机磷酸酯、拟除虫菊酯、百草枯
    \item \imp{其他毒物}：氯乙烯、含氟塑料、丙烯腈等
\end{enumerate}
\end{keybox}

% ----- 3.2 毒物的体内过程 -----
\section{生产性毒物的体内过程}

\subsection{毒物的存在状态与接触途径}

\begin{defbox}
生产性毒物主要以\keyword{气溶胶}和\keyword{气态}两大类形式存在：

\textbf{气溶胶：}
\begin{itemize}[leftmargin=*]
    \item 液态——雾（Mist）：悬浮于空气中的液态小微粒
    \item 固态——粉尘（Dust）：悬浮于空气中粒径0.1--10$\mu$m的固态小微粒
    \item 固态——烟（Smoke）：悬浮于空气中粒径<0.1$\mu$m的固态小微粒
\end{itemize}

\textbf{气态：}
\begin{itemize}[leftmargin=*]
    \item 气体：常温常压下呈气态的物质
    \item 蒸气：固体升华或液体蒸发形成
\end{itemize}
\end{defbox}

\subsection{毒物的吸收途径}

\begin{keybox}
\textbf{呼吸道（最主要途径）：}
\begin{itemize}[leftmargin=*]
    \item 气态毒物：受空气中浓度/分压、分子量、血/气分配系数、水溶性、劳动强度、肺通气量、肺血流量及气象条件等影响
    \item 气溶胶毒物：受气道结构、粒子形状、粒子分散度、溶解度、呼吸道清除功能影响
\end{itemize}

\textbf{消化道：}
\begin{itemize}[leftmargin=*]
    \item 一般通过简单扩散吸收；胃液pH极低（2.0），弱有机酸多未解离易吸收
    \item 高脂溶性物质（有机氯农药等）可经肠道淋巴系统吸收
\end{itemize}

\textbf{皮肤：}
\begin{itemize}[leftmargin=*]
    \item 脂溶性物质可经皮肤吸收；某些物质（如有机磷农药）经皮肤吸收是重要的中毒途径
\end{itemize}
\end{keybox}

\subsection{毒物在体内的分布与蓄积}

\begin{defbox}
\textbf{贮存库（Storage Depot）}：毒物在体内的蓄积部位，与血浆中游离毒物保持动态平衡。
\begin{itemize}[leftmargin=*]
    \item 血浆贮存库——血浆蛋白结合
    \item 肝贮存库——金属硫蛋白等
    \item 肾贮存库——金属硫蛋白
    \item 脂肪组织贮存库——脂溶性毒物（有机氯农药等）
    \item 骨骼贮存库——铅、氟等（占体内铅90\%）
\end{itemize}
\end{defbox}

\subsection{毒物的排泄途径}

\begin{enumerate}[leftmargin=*]
    \item \imp{经尿液排泄}：肾小球滤过+肾小管简单扩散+主动转运。脂溶性物质可被重吸收，只有水溶性/离子型物质进入尿液
    \item \imp{经胆汁排泄}：肝脏代谢产物排入胆汁→粪便排出。部分经\keyword{肝肠循环}重吸收
    \item \imp{经肺呼出}：气态化合物经呼吸道排出，速度取决于血中溶解度、呼吸速度、肺血流速度
    \item \imp{其他途径}：汗液、唾液、乳汁（有机氯杀虫剂、多卤联苯、咖啡碱、某些金属可随乳汁排出，对婴儿有损害意义）
\end{enumerate}

\subsection{毒物的生物转化}

\begin{keybox}
\textbf{生物转化的意义：}
\begin{itemize}[leftmargin=*]
    \item 使污染物极性和水溶性增高，易于排出
    \item 使一些污染物生物学活性降低或消除（灭活/解毒作用）
    \item 但也可使某些物质毒性增强（\keyword{代谢活化}）：
    \begin{itemize}
        \item 代谢解毒：化学物（毒性）→中间产物（低毒/无毒）→产物（无毒）
        \item \textbf{代谢活化}：化学物（无毒）→中间产物（毒性）→产物（无毒）
    \end{itemize}
\end{itemize}

\textbf{生物转化反应分类：}
\begin{enumerate}[leftmargin=*]
    \item \imp{I相代谢}：氧化、还原、水解——官能团化反应，引入极性基团（-OH、-COOH、-NH$_2$、-SH）
    \item \imp{II相代谢（结合反应）}：与葡萄糖醛酸、硫酸等内源性小分子结合，产生极性强或水溶性的结合物排出
\end{enumerate}

\textbf{发生部位：}肝（主要）、肾、胃肠道、肺、皮肤、胎盘。肝细胞定位：微粒体、胞液、线粒体。

\textbf{影响生物转化的因素：}
\begin{itemize}[leftmargin=*]
    \item 物种差异和个体差异（如苯胺在小鼠t$_{1/2}$=35min，狗t$_{1/2}$=167min）
    \item 毒物代谢酶的诱导和激活、抑制和阻遏
    \item 性别、遗传、年龄、营养、疾病等
\end{itemize}
\end{keybox}

% ----- 3.3 职业中毒的临床 -----
\section{职业中毒的临床类型与诊断}

\subsection{临床类型}

\begin{table}[H]
\centering
\caption{职业中毒的临床类型}
\begin{tabularx}{\textwidth}{llX}
\hline
\textbf{类型} & \textbf{接触特征} & \textbf{典型举例} \\
\hline
\imp{急性中毒} & 一次或短时间（几分钟至数小时）内大量接触 & 急性苯中毒、氯气中毒 \\
\imp{慢性中毒} & 少量、长期接触 & 慢性铅中毒、锰中毒 \\
亚急性中毒 & 介于急性和慢性之间 & 亚急性铅中毒 \\
迟发性中毒 & 脱离接触毒物一定时间后发病 & 迟发性锰中毒 \\
\hline
\end{tabularx}
\end{table}

\subsection{不同靶系统的中毒表现}

\begin{table}[H]
\centering
\caption{职业中毒的靶系统及临床表现}
\begin{tabularx}{\textwidth}{lXl}
\hline
\textbf{靶系统} & \textbf{临床症状} & \textbf{典型毒物} \\
\hline
\imp{神经系统} & 周围神经病变、锥体外系症状、中毒性脑病 & 重金属、有机溶剂、窒息性气体 \\
\imp{呼吸系统} & 呼吸道炎症、化学性肺炎、肺水肿 & 刺激性气体、窒息性气体 \\
\imp{血液系统} & 造血功能抑制、血红蛋白变性 & 重金属、苯系化合物 \\
\imp{消化系统} & 中毒性肝病、腹痛 & 重金属、农药、有机溶剂 \\
\imp{泌尿系统} & 中毒性肾病、泌尿系统肿瘤 & 重金属、苯系化合物 \\
\imp{循环系统} & 心肌损害、心律失常 & 重金属、有机溶剂 \\
皮肤 & 接触性皮炎、职业性痤疮 & 酸碱、有机溶剂、石油产品 \\
\hline
\end{tabularx}
\end{table}

\subsection{职业中毒的诊断与治疗}

\begin{keybox}
\textbf{诊断依据：}
\begin{enumerate}[leftmargin=*]
    \item 详细的职业接触史
    \item 相应的临床表现和辅助检查结果
    \item 职业卫生现场调查（车间空气测定等）
    \item 综合分析，排除其他病因所致类似疾病
\end{enumerate}

\textbf{治疗原则：}
\begin{enumerate}[leftmargin=*]
    \item \imp{病因治疗}：使用特效解毒剂（络合剂、高铁血红蛋白生成剂等）
    \item \imp{对症治疗}：缓解临床症状
    \item \imp{支持疗法}：维持生命体征，促进康复
\end{enumerate}

\textbf{预防原则：}
\begin{enumerate}[leftmargin=*]
    \item 无毒/低毒材料替代
    \item 降低车间空气中毒物浓度（工艺改革、通风）
    \item 加强个人防护（防护服、防毒口罩）
    \item 定期监测和职业健康检查
\end{enumerate}
\end{keybox}

% ----- 3.4 金属与类金属中毒 -----
\section{金属与类金属中毒}

\subsection{金属中毒概述}

\begin{defbox}
\textbf{必需微量元素}：铜、铁、锌、硒、铬、碘、钴、钼
\textbf{可能必需微量元素}：锰、硅、镍、钒
\textbf{潜在毒性微量元素}：氟、汞、砷、铝、锡、锂、铅（低剂量可能对人体有作用）

金属对人体的影响机制：
\begin{enumerate}[leftmargin=*]
    \item 与\keyword{巯基（-SH）}共价结合，改变生物大分子结构和功能
    \item 产生过多自由基，破坏氧化还原系统，引起氧化损伤
    \item 与必需微量元素相互作用，干扰正常生理生化功能
    \item 诱导细胞合成保护性蛋白
    \item 表观遗传调控（DNA甲基化、非编码RNA、组蛋白修饰）
\end{enumerate}

\textbf{络合剂治疗：}
\begin{itemize}[leftmargin=*]
    \item \imp{氨羧络合剂}：如依地酸二钠钙（CaNa$_2$-EDTA），与多种金属络合成无毒络合物排出
    \item \imp{巯基络合剂}：如二巯基丙醇、二巯基丙磺酸钠、二巯基丁二酸钠、青霉胺，保护巯基酶系统
\end{itemize}
\end{defbox}

% ----- 3.4.1 铅中毒 -----
\subsection{铅中毒（Lead Poisoning）}

\begin{keybox}
\textbf{理化特性：}柔软、略带灰白色重金属，熔点327$^\circ$C；加热至400--500$^\circ$C有大量铅蒸气逸出，在空气中迅速氧化为PbO。

\textbf{接触机会：}
\begin{itemize}[leftmargin=*]
    \item 职业暴露（呼吸道吸入）：铅矿开采冶炼、含铅蓄电池生产回收、含铅制品制造
    \item 生活暴露（消化道）：含铅输水管道、含铅颜料、草药（铅丹Pb$_3$O$_4$、密陀僧PbO）
\end{itemize}

\textbf{毒作用机制——卟啉代谢障碍（血液系统）*：}

铅抑制血红素合成过程中的关键酶：
\begin{enumerate}[leftmargin=*]
    \item 抑制\keyword{$\delta$-氨基乙酰丙酸脱水酶（ALAD）}→ 血/尿ALA↑
    \item 抑制\keyword{亚铁螯合酶}→ 红细胞游离原卟啉（FEP）↑、锌原卟啉（ZPP）↑
    \item 血红素合成障碍 → 血红蛋白↓ → 低色素性贫血
    \item 嗜碱性点彩红细胞↑、网织红细胞↑
\end{enumerate}

\textbf{铅中毒的异常血红素前体：}ALA↑（血、尿）；FEP↑、ZPP↑（红细胞）

\textbf{其他毒作用机制：}
\begin{itemize}[leftmargin=*]
    \item \imp{神经系统}：ALA升高→GABA受体竞争性抑制→神经效应改变；损害周围神经→脱髓鞘→神经传导速度减慢→\textbf{垂腕、垂足}
    \item \imp{泌尿系统}：损害近曲小管→Fanconi综合征（尿糖、尿氨基酸、尿磷酸盐↑）
    \item \imp{消化系统}：抑制肠壁碱性磷酸酶和ATP酶→肠壁平滑肌痉挛→\textbf{铅绞痛}
\end{itemize}
\end{keybox}

\begin{keybox}
\textbf{慢性铅中毒临床表现（较常见）：}
\begin{itemize}[leftmargin=*]
    \item \imp{神经系统}：
    \begin{itemize}
        \item 类神经症：头昏、头痛、无力、记忆力减退、睡眠障碍（功能性）
        \item 周围神经病：感觉型（肢端麻木、手套袜套型感觉障碍）；运动型（肌无力、腕下垂、足下垂）；混合型
        \item 中毒性脑病（最严重，少见）：头痛、恶心、呕吐、高热、烦躁、抽搐、昏迷
    \end{itemize}
    \item \imp{造血系统}：低色素正常细胞/小细胞性贫血（血浆铁不降低）；点彩红细胞↑、网织红细胞↑
    \item \imp{消化系统}：食欲不振、恶心、铅线（齿龈边缘蓝黑色PbS沉积）、\keyword{铅绞痛}（突发剧烈腹痛）
    \item \imp{其他}：铅容（面色苍白）、高血压、铅性肾病
\end{itemize}

\textbf{急性铅中毒（较少见）：}经口服用大量铅化合物。潜伏期4--6小时至1--2周。铅绞痛、口内金属味、铅容。

\textbf{诊断（GBZ37-2024）：}职业接触史（密切接触铅$\geq$3个月）+车间空气测定+临床症状（神经、消化、造血系统）+实验室检查（血铅、尿铅、$\delta$-ALA、FEP、ZPP）

\textbf{治疗：}金属络合剂驱铅治疗——依地酸二钠钙（CaNa$_2$-EDTA）；对症治疗。血铅>700$\mu$g/L尤其有中枢/周围神经紊乱时考虑络合剂驱铅。

\textbf{预防：}无毒/低毒材料替代（锌钡白代替铅白；铁红代替铅丹；激光排版代替铅字）；降低车间铅浓度（工艺改革、通风，控制熔铅温度减少铅蒸气）；个人防护（防护服、防烟尘口罩）；定期健康检查（妊娠及哺乳期女工调离）
\end{keybox}

% ----- 3.4.2 汞中毒 -----
\subsection{汞中毒（Mercury Poisoning）}

\begin{keybox}
\textbf{接触机会：}汞矿开采冶炼、温度计/气压计/血压计制造（逐步替代）、催化剂、汞齐生产应用、军工雷汞、核工业冷却剂、含汞化妆品、误服汞盐

\textbf{体内过程：}
\begin{itemize}[leftmargin=*]
    \item 吸收：呼吸道（金属汞约85\%经肺泡壁吸收）；消化道（汞盐及有机汞，>0.01\%）；皮肤基本不吸收
    \item 分布不均：肾中最高，其次为肝、脑
    \item 蓄积：肾（金属硫蛋白）
    \item 排泄：以尿汞为主
\end{itemize}

\textbf{慢性汞中毒临床表现*（较常见，6个月及以上发病）：}
\begin{itemize}[leftmargin=*]
    \item \imp{易兴奋症}：急躁、易怒、胆怯、性格改变等——\textbf{特有表现}
    \item \imp{震颤}：眼睑、舌尖、手指→意向性大震颤（动作开始时震颤，动作过程中加重，动作完成后停止；紧张或刻意控制时明显加重）
    \item \imp{口腔炎}：流涎、汞线、牙龈肿胀
\end{itemize}

\textbf{急性汞中毒：}短时间吸入高浓度汞蒸气或摄入可溶性汞盐。起病急骤，发热、咳嗽、呼吸困难、口腔炎和胃肠道症状，可发生化学性肺炎、汞毒性皮炎。口服升汞→急性口腔炎/腐蚀性肠胃炎、汞毒性肾炎、蛋白尿甚至肾衰竭。

\textbf{诊断（GBZ 89-2024）：}金属汞职业接触史+临床表现+职业卫生学调查+综合分析。尿汞生物接触限值：20$\mu$mol/mol肌酐。驱汞试验：5\%二巯丙磺钠5mL肌注，尿汞>45$\mu$g/L提示过量汞吸收。

\textbf{治疗：}
\begin{itemize}[leftmargin=*]
    \item 误服汞盐：\textbf{不能洗胃}，应灌服高蛋白物质（蛋清/牛奶/豆浆）+活性炭吸附+硫酸镁导泻
    \item 驱汞治疗：二巯基丙磺酸钠、二巯基丁二酸、二巯基丁二酸钠
\end{itemize}

\textbf{汞中毒的独特预防措施（与其他金属中毒不同）：}
\begin{enumerate}[leftmargin=*]
    \item 车间地面、工作台面应光滑无裂缝，便于冲洗（防止汞滴积聚）
    \item 墙壁和天花板应涂刷过氯乙烯漆（减少汞蒸气吸附）
    \item 定期用1\%碘蒸气熏蒸车间（碘与汞生成不挥发的碘化汞）
    \item 作业场所温度控制在15--16$^\circ$C以下（降低汞蒸发速度）
    \item 加强通风，在汞发生源设置槽边吸风装置
    \item 定期进行尿汞检测和职业健康检查
\end{enumerate}
\end{keybox}

% ----- 3.4.3 锰中毒 -----
\subsection{锰中毒（Manganese Poisoning）}

\begin{keybox}
\textbf{接触机会：}锰矿开采冶炼、干电池制造（二氧化锰为原料）、电焊条制造及使用、染料工业

\textbf{中毒机制（尚未完全清楚*）：}
\begin{enumerate}[leftmargin=*]
    \item 使中枢神经突触线粒体ATP酶活性降低 → 影响神经突触传导能力
    \item 抑制胆碱酯酶合成 → 乙酰胆碱蓄积 → \textbf{肌张力增高}（震颤麻痹）
    \item 基底神经节、纹状体内\keyword{5-羟色胺和多巴胺含量降低} → 神经传导障碍
\end{enumerate}

\textbf{慢性锰中毒临床表现（较多见）*：}
\begin{itemize}[leftmargin=*]
    \item 早期：类神经症和自主神经功能障碍（记忆力减退、嗜睡、精神萎靡）
    \item 继而：锥体外系神经受损——肌张力增高、手指细小震颤、腱反射亢进、激动多汗
    \item 后期：\keyword{帕金森氏综合征}——说话含糊不清、面部表情减少（面具脸）、动作笨拙、慌张步态、肌张力齿轮样增强、静止性震颤、记忆力减退、智能下降、强迫观念
\end{itemize}

\textbf{急性锰中毒（较少见）：}口服高锰酸钾→腐蚀性胃肠炎/刺激性支气管炎；电焊→金属烟热（咽痛、咳嗽、寒颤高热）

\textbf{治疗：}
\begin{itemize}[leftmargin=*]
    \item 驱锰：依地酸二钠钙（CaNa$_2$-EDTA）、二巯基丁二酸、对氨基水杨酸钠（PAS）
    \item 对症（震颤麻痹）：左旋多巴制剂、金刚烷胺
    \item 急救（口服KMnO$_4$）：温水洗胃、口服牛奶和氢氧化铝凝胶
\end{itemize}

\textbf{预防：}通风降尘、焊接密闭吸尘装置、自动电焊代替手工、佩戴滤膜口罩
\end{keybox}

\subsection{铅、汞、锰中毒对比总结}

\begin{table}[H]
\centering
\caption{铅、汞、锰中毒对比}
\begin{tabularx}{\textwidth}{lXXX}
\hline
\textbf{特征} & \textbf{铅（Pb）} & \textbf{汞（Hg）} & \textbf{锰（Mn）} \\
\hline
主要吸收途径 & 呼吸道 & 呼吸道（85\%） & 呼吸道 \\
主要排泄 & 尿液（血铅、尿铅） & 尿液（尿汞） & 胆道→粪便 \\
毒作用机制 & 卟啉代谢障碍、巯基结合 & 与巯基结合、抑制酶活性 & 线粒体ATP酶↓、ACh蓄积、多巴胺↓ \\
特征性表现 & 铅绞痛、垂腕、铅线、贫血 & 易兴奋症、意向性震颤、口腔炎 & 帕金森综合征、面具脸、齿轮样肌张力 \\
解毒剂 & CaNa$_2$-EDTA & 二巯基丙磺酸钠 & CaNa$_2$-EDTA、PAS \\
\hline
\end{tabularx}
\end{table}

% ----- 3.5 刺激性气体 -----
\section{刺激性气体}

\begin{defbox}
\textbf{刺激性气体（Irritant Gases）}：对眼、呼吸道黏膜及皮肤有刺激作用，引起机体以急性炎症、\keyword{肺水肿}为主要病理改变的一类气态物质。多具腐蚀性。

\textbf{常见类型：}Cl$_2$、NH$_3$、NOx、SO$_2$、COCl$_2$（光气）、SO$_3$、甲醛

\textbf{分类：}酸类（硫酸、硝酸、盐酸）；氮的氧化物；硫的化合物；氯及其化合物；醛类；酯类；卤烃类
\end{defbox}

\subsection{刺激性气体的毒作用机制与特点}

\begin{keybox}
\textbf{毒作用机制：}
\begin{enumerate}[leftmargin=*]
    \item \imp{直接刺激腐蚀作用}：酸——吸收组织水分、凝固蛋白、细胞坏死；碱——吸收细胞水分、皂化脂肪、细胞坏死
    \item \imp{引起呼吸道炎症反应}：刺激细胞因子"瀑布效应"扩大炎症损伤
    \item \imp{自由基损伤作用}：自身或产物为自由基，基于炎症反应的自由基
\end{enumerate}

\textbf{毒作用特点*：}
\begin{itemize}[leftmargin=*]
    \item 病变程度主要取决于\keyword{吸入浓度和持续接触时间}
    \item 病变部位与\keyword{毒物水溶性}密切相关：
    \begin{itemize}
        \item \textbf{水溶性大}（NH$_3$、HCl）：主要作用于上呼吸道，刺激症状明显
        \item \textbf{水溶性中等}（Cl$_2$、SO$_2$）：作用于整个呼吸道
        \item \textbf{水溶性小}（光气、NO$_2$）：主要作用于深呼吸道和肺泡，刺激症状轻但危害重
    \end{itemize}
    \item 接触时间短-浓度低：局部刺激为主
    \item 浓度高/时间长：局部损害和全身反应（喉痉挛、支气管痉挛；反射性中枢抑制、昏迷或休克）
\end{itemize}
\end{keybox}

\subsection{中毒性肺水肿}

\begin{keybox}
\textbf{中毒性肺水肿（Toxic Pulmonary Edema）*}：吸入高浓度刺激性气体后引起肺泡及肺间质过量体液潴留的病理过程，最终可导致急性呼吸功能衰竭。是刺激性气体所致\keyword{最严重危害}和\keyword{职业病最常见急症}之一。

\textbf{常见引起肺水肿的气体：}光气、二氧化氮、氨、氯、臭氧、硫酸二甲酯、羰基镍

\textbf{发病机制：}刺激性气体 → 神经-体液反射/局部刺激 → 肺血管、淋巴管痉挛 → 肺淋巴循环梗阻 → 血管活性物质释放 → 肺泡表面活性物质减少 → 肺泡/毛细血管通透性增加 → 肺水肿 → 缺氧 → 休克 → ARDS

\textbf{临床过程（四期）*：}
\begin{enumerate}[leftmargin=*]
    \item \imp{刺激期}：气管-支气管黏膜急性炎症。水溶性低的气体此期症状不明显
    \item \imp{潜伏期}（一般2--6小时）：\textbf{假象期}——症状减轻或消失，无明显体征。\textbf{病情转归的关键期}
    \item \imp{肺水肿期}（24h内剧变）：症状突然加重——呼吸困难、烦躁、大汗淋漓、咳大量粉红色泡沫痰、口唇发绀、两肺湿啰音。胸片可见蝴蝶形阴影，持续1--3天
    \item \imp{恢复期}（2--15天）：1--2周内逐渐恢复，多无后遗症
\end{enumerate}

\textbf{救治原则（关键是积极防治肺水肿和ARDS）：}
\begin{enumerate}[leftmargin=*]
    \item 现场急救处理（脱离接触、清洗皮肤眼）
    \item 迅速纠正缺氧（氧疗浓度50\%，机械通气）
    \item 糖皮质激素（\textbf{早期、足量、短程}使用，降低肺毛细血管通透性）
    \item 去泡沫剂二甲硅酮、维持呼吸道通畅
    \item 控制液体入量、预防继发感染
\end{enumerate}
\end{keybox}

\subsection{氯气（Cl$_2$）}

\begin{keybox}
\textbf{理化特性：}黄绿色、有异臭和强烈刺激性气体，易溶于水和碱性溶液。遇水生成次氯酸和盐酸；高温与CO作用生成光气。

\textbf{接触机会：}电解食盐、制造氯化物（四氯化碳、漂白粉、环氧树脂）、制药/皮革/造纸/印染业作氧化剂或漂白剂

\textbf{急性中毒分级：}
\begin{itemize}[leftmargin=*]
    \item 刺激反应：眼、呼吸道刺激
    \item 轻度：支气管炎或支气管周围炎
    \item 中度：支气管肺炎、间质性/局限性肺泡性肺水肿
    \item 重度：弥漫性肺泡性肺水肿或中央性肺水肿
\end{itemize}

\textbf{慢性影响：}慢性支气管炎、支气管哮喘、肺气肿；类神经症；皮肤痤疮样皮疹或疱疹

\textbf{治疗：}现场处理→合理氧疗→糖皮质激素（早期、足量、短程）→维持呼吸道通畅→控制液体入量→预防继发感染。皮肤接触液氯→灼伤或急性皮炎。

\textbf{预防：}防止跑冒滴漏、保持管道负压、加强通风和密闭操作、含氯废气石灰净化处理、MAC=1 mg/m$^3$
\end{keybox}

% ----- 3.6 窒息性气体 -----
\section{窒息性气体}

\begin{defbox}
\textbf{窒息性气体的分类：}
\begin{enumerate}[leftmargin=*]
    \item \imp{单纯窒息性气体}：无毒或惰性气体，使空气中氧含量降低 → 机体缺氧（N$_2$、CH$_4$、CO$_2$）
    \item \imp{化学窒息性气体}：
    \begin{itemize}
        \item 血液窒息性：使血液运氧能力下降（如CO竞争性结合Hb）
        \item 细胞窒息性：使组织利用氧的能力障碍（如HCN、H$_2$S抑制细胞呼吸酶）
    \end{itemize}
\end{enumerate}
\end{defbox}

\begin{keybox}
\textbf{毒作用特点*：}
\begin{itemize}[leftmargin=*]
    \item 主要致病环节均为\keyword{引起机体组织细胞缺氧}
    \item \keyword{脑对缺氧极为敏感}——脑重量仅占体重2--3\%，但耗氧量占总耗氧量20--30\%，无供能物质储备，全靠脑循环供给氧气
    \item \textbf{防治脑水肿是治疗的关键}
    \item 慢性中毒尚未确证
\end{itemize}

\textbf{临床表现：}
\begin{itemize}[leftmargin=*]
    \item 缺氧表现：中枢神经系统（头痛、烦躁、肌肉抽搐、语言障碍、嗜睡、昏迷）；呼吸循环系统（呼吸浅促、紫绀、心跳过速、血压下降、休克）
    \item 脑水肿（严重缺氧或窒息3--5min）：颅压增高、头痛、呕吐、血压升高、心率减慢
    \item 面部颜色：CO中毒呈\keyword{樱桃红色}；氰化物中毒可有心肌损害、肺水肿
\end{itemize}

\textbf{治疗原则（防治脑水肿是关键）：}
\begin{enumerate}[leftmargin=*]
    \item 积极纠正脑缺氧（吸氧40--60\%）
    \item 解毒治疗；降低颅压和解除脑水肿（肾上腺糖皮质激素、脱水剂）
    \item 降低血液粘稠度；改善脑组织代谢
    \item 清除自由基；钙通道阻滞剂；对症及支持治疗
\end{enumerate}
\end{keybox}

\subsection{一氧化碳（CO）中毒}

\begin{keybox}
\textbf{理化特性：}无色、无臭、无味、无刺激；易燃易爆；微溶于水；不易被活性炭吸附。\textbf{CO中毒是我国发病和死亡人数最高的急性职业中毒。}

\textbf{中毒机制*：}
\begin{enumerate}[leftmargin=*]
    \item CO与Hb（Fe$^{2+}$）结合 → \keyword{碳氧血红蛋白（HbCO）}——\textbf{亲和力比O$_2$大200--300倍}，解离速度比O$_2$慢3600倍 → 阻碍氧运输
    \item 与肌红蛋白结合 → 损害氧的细胞弥散和线粒体功能
    \item 与还原型细胞色素氧化酶结合 → 阻碍氧利用
\end{enumerate}

\textbf{急性CO中毒分级（依据HbCO\%）：}
\begin{itemize}[leftmargin=*]
    \item \imp{轻度}（HbCO >10\%）：剧烈头痛、头晕、四肢无力、恶心呕吐、无昏迷
    \item \imp{中度}（HbCO >30\%）：面色潮红、多汗、脉快、浅至中度昏迷
    \item \imp{重度}（HbCO >50\%）：深度昏迷、脑水肿、休克、严重心肌损害、肺水肿、呼吸衰竭
    \item \imp{迟发性脑病*}：\textbf{2--60天"假愈期"}后出现精神及意识障碍（痴呆状态）、锥体外系/锥体系神经损害
\end{itemize}

\textbf{实验室检查：}血中HbCO测定（t$_{1/2}$=5h）；脑电图及诱发电位；脑CT与核磁共振；心肌酶检查；心电图

\textbf{治疗：}立即脱离中毒环境→高浓度氧疗（高压氧治疗）→防治脑水肿→对症支持
\end{keybox}

\subsection{氰化氢（HCN）中毒}

\begin{keybox}
\textbf{理化特性：}无色、具\keyword{苦杏仁味}气体，密度0.93，易扩散，易溶于水、脂肪及有机溶剂。

\textbf{接触机会：}氰化物生产、电镀、贵重金属提炼、制药、合成材料、灭鼠剂/杀虫剂生产、钢铁热处理。\textbf{日常生活：}苦杏仁、桃仁、枇杷仁、木薯食用不当中毒。

\textbf{中毒机制*：}CN$^-$与细胞色素氧化酶Fe$^{3+}$结合 → 氰化高铁细胞色素氧化酶 → Fe$^{3+}$失去传递电子能力 → \keyword{中断呼吸链}→ 阻断细胞内氧化代谢 → \keyword{细胞窒息死亡}。氰化物可经呼吸道、皮肤、消化道进入机体。大部分在硫氰酸酶作用下转化为硫氰酸盐随尿排出。

\textbf{急性中毒四期（未发生"电击样"死亡者）：}
\begin{enumerate}[leftmargin=*]
    \item \imp{前驱期}：刺激症状、呼气带苦杏仁味
    \item \imp{呼吸困难期}：胸闷、心悸、呼吸困难
    \item \imp{痉挛期}：强直性痉挛
    \item \imp{麻痹期}：肌肉松弛、反射消失，重者昏迷死亡
\end{enumerate}

\textbf{特效解毒剂——"亚硝酸钠-硫代硫酸钠"疗法：}
\begin{enumerate}[leftmargin=*]
    \item 高铁血红蛋白生成剂（亚硝酸钠/4-DMAP）将Fe$^{2+}$-Hb氧化为Fe$^{3+}$-Hb → 后者与CN$^-$结合 → 氰化高铁血红蛋白（暂时解毒）
    \item 硫代硫酸钠在硫氰酸酶/转硫酶作用下 → 将CN$^-$转化为无毒硫氰酸盐（NaSCN）→ 随尿排出
\end{enumerate}
\end{keybox}

\subsection{硫化氢（H$_2$S）}

H$_2$S是具有臭蛋味的无色气体，比重1.19。主要由含硫有机物腐败产生，常见于下水道清理、粪坑、腌渍池、造纸厂等场所。中毒机制与HCN类似，抑制细胞色素氧化酶。治疗可用高铁血红蛋白生成剂。

% ----- 3.7 有机溶剂 -----
\section{有机溶剂}

\subsection{概述}

有机溶剂多为脂溶性、易挥发的有机液体，主要经呼吸道吸入，也可经皮肤吸收。共同毒性特点：\textbf{中枢神经系统抑制作用（麻醉作用）}；对皮肤黏膜有刺激作用。

\subsection{苯（Benzene）}

\begin{keybox}
\textbf{接触机会：}苯的生产（煤焦油分馏、石油裂解）、用作溶剂和稀释剂（油漆、喷漆、制鞋、箱包）、化工原料（合成橡胶、塑料等）

\textbf{毒作用机制*：}
\begin{itemize}[leftmargin=*]
    \item 代谢产物（酚类、氢醌、苯醌等）在骨髓中蓄积 → \keyword{抑制造血功能}
    \item 损害骨髓造血干细胞和骨髓基质细胞
    \item 致染色体畸变和白血病
\end{itemize}

\textbf{临床表现：}
\begin{itemize}[leftmargin=*]
    \item \imp{急性中毒}：中枢神经系统麻醉症状——头痛、眩晕、步态蹒跚、醉酒感；重者昏迷、抽搐、呼吸循环衰竭
    \item \imp{慢性中毒}：
    \begin{itemize}
        \item 造血系统损害（最特征性）：\textbf{白细胞减少→血小板减少→全血细胞减少→再生障碍性贫血→白血病}
        \item 神经系统：类神经症
    \end{itemize}
\end{itemize}

\textbf{诊断（GBZ 68-2022）：}职业接触史+血常规异常+排除其他原因。苯为\keyword{IARC 1类致癌物}（致白血病）。

\textbf{预防：}以无毒/低毒溶剂替代（甲苯、二甲苯代替苯作溶剂）；加强通风排毒；个人防护。
\end{keybox}

\subsection{其他常见有机溶剂}

\begin{table}[H]
\centering
\caption{常见有机溶剂的中毒特点}
\begin{tabularx}{\textwidth}{lX}
\hline
\textbf{名称} & \textbf{主要毒作用及临床特点} \\
\hline
\imp{甲苯、二甲苯} & 麻醉作用较苯强；对造血系统损害较苯弱；皮肤黏膜刺激 \\
\imp{正己烷} & 主要损害\keyword{周围神经系统}——远端轴索病、手套袜套样感觉障碍、肌无力 \\
\imp{二硫化碳（CS$_2$）} & 神经系统损害（类神经症、周围神经病）；心血管系统（动脉粥样硬化）\\
\imp{四氯化碳（CCl$_4$）} & \keyword{肝毒性}（中毒性肝炎）和肾毒性；麻醉作用 \\
\hline
\end{tabularx}
\end{table}

% ----- 3.8 苯的氨基和硝基化合物 -----
\section{苯的氨基和硝基化合物}

\begin{keybox}
\textbf{共同毒作用特点：}
\begin{enumerate}[leftmargin=*]
    \item \imp{形成高铁血红蛋白（MetHb）}：使Fe$^{2+}$-Hb氧化为Fe$^{3+}$-Hb → 失去携氧能力 → 组织缺氧（\keyword{化学性青紫}）
    \item \imp{溶血作用}：赫恩氏小体（Heinz body）形成 → 溶血性贫血
    \item \imp{肝脏损害}：中毒性肝炎
    \item \imp{泌尿系统损害}：出血性膀胱炎
    \item \imp{致癌作用}：苯胺致膀胱癌（IARC 1类）；联苯胺致膀胱癌
\end{enumerate}

\begin{exambox}
\textbf{常见物质形成高铁血红蛋白能力的强弱比较（考试常考排序）：}\\
对硝基苯 $>$ 间位二硝基苯 $>$ 苯胺 $>$ 邻位二硝基苯 $>$ 硝基苯

例外：二硝基酚、联苯胺等不形成高铁血红蛋白。
\end{exambox}

\textbf{高铁血红蛋白血症的治疗：}亚甲蓝（美蓝）——将Fe$^{3+}$-Hb还原为Fe$^{2+}$-Hb。需小剂量使用（1--2mg/kg），大剂量反而氧化Hb。

\subsection{苯胺（Aniline）}
经呼吸道、皮肤和消化道吸收，经皮肤吸收易被忽视而成为引起职业中毒的主因。代谢活化生成苯基羟胺（苯胲），进一步转化为对氨基酚。主要表现为：
\begin{itemize}[leftmargin=*]
    \item \imp{高铁血红蛋白血症（化学性发绀，蓝灰色）}：
    \begin{itemize}
        \item MetHb 15\%：明显发绀，无明显自觉症状
        \item MetHb 30\%：头昏/头痛/乏力/恶心/手指麻木/视物模糊
        \item MetHb 50\%：中枢神经系统缺氧症状；休克，瞳孔散大，反应消失
    \end{itemize}
    \item \imp{溶血性贫血}：赫恩滋小体（Heinz body）形成
    \item \imp{中毒性肝炎}和\imp{膀胱刺激症状}
    \item 尿中对氨基酚增高（接触指标）
\end{itemize}

\subsection{三硝基甲苯（TNT）}
经呼吸道、皮肤和消化道吸收，TNT具有亲脂性，很容易从皮肤吸收。主要损害：
\begin{itemize}[leftmargin=*]
    \item \keyword{中毒性白内障}（最特征性改变，晶状体混浊*）：6月--3年可发生，工龄越长发病率越高（10年78.5\%，15年83.65\%）；白内障形成后即使不再接触TNT仍可进展；一般不影响视力，但晶体中央部混浊可使视力下降；与中毒性肝病发病不平行
    \item \keyword{中毒性肝炎}（中毒性肝病）：肝细胞坏死、脂肪变性、中毒性肝硬化
    \item \keyword{再生障碍性贫血}：赫恩小体形成
    \item 其他：睾酮降低、精子形成受损、女性月经异常
    \item 生物监测指标：尿中4-A（4-氨基-2,6-二硝基甲苯）和原形TNT
\end{itemize}
\end{keybox}

% ----- 3.9 高分子化合物 -----
\section{高分子化合物生产中的毒物}

\begin{table}[H]
\centering
\caption{高分子化合物主要职业中毒}
\begin{tabularx}{\textwidth}{lX}
\hline
\textbf{名称} & \textbf{接触机会与主要毒作用} \\
\hline
\imp{氯乙烯} & 聚氯乙烯（PVC）生产；\keyword{肝血管肉瘤}（IARC 1类致癌物*）；肢端溶骨症；肝纤维化 \\
\imp{含氟塑料热解物} & 聚四氟乙烯（PTFE）高温热解产生；\keyword{聚合物烟雾热}；肺纤维化 \\
\imp{丙烯腈} & 合成腈纶纤维的原料；类似氰化物中毒机制（释放CN$^-$）；高铁血红蛋白血症 \\
\hline
\end{tabularx}
\end{table}

% ----- 3.10 农药中毒 -----
\section{农药中毒}

\subsection{有机磷酸酯类农药}

\begin{keybox}
\textbf{毒作用机制*：}不可逆性抑制\keyword{乙酰胆碱酯酶（AChE）}→ 乙酰胆碱（ACh）蓄积 → 胆碱能神经过度兴奋 → 毒蕈碱样症状、烟碱样症状和中枢神经系统症状。

\textbf{临床表现：}
\begin{itemize}[leftmargin=*]
    \item \imp{毒蕈碱样症状（M样）}：瞳孔缩小、视力模糊、流泪流涎、恶心呕吐、腹痛腹泻、大小便失禁、支气管痉挛、呼吸困难、肺水肿
    \item \imp{烟碱样症状（N样）}：肌束震颤、肌无力、肌麻痹（呼吸肌麻痹致死亡）
    \item \imp{中枢神经系统症状}：头痛、头晕、抽搐、昏迷
\end{itemize}

\textbf{诊断依据：}接触史+M样和N样症状+全血ChE活性降低（<70\%轻度，<50\%中度，<30\%重度）

\textbf{治疗：}
\begin{enumerate}[leftmargin=*]
    \item 清除毒物（脱离接触、洗胃、皮肤清洗）
    \item \imp{特效解毒剂——阿托品}（对抗M样症状，\textbf{早期、足量、反复给药}，阿托品化标志：瞳孔扩大、颜面潮红、口干、心率加快）
    \item \imp{胆碱酯酶复能剂——氯磷定/解磷定}（恢复ChE活性，对N样症状有效，对老化的ChE无效）
\end{enumerate}

\textbf{中间期肌无力综合征（IMS）：}急性中毒后1--4天，出现颈屈肌、四肢近端肌和呼吸肌无力。

\textbf{迟发性多神经病（OPIDN）：}急性中毒后2--3周出现，四肢远端感觉和运动障碍。
\end{keybox}

\subsection{其他农药}

\begin{itemize}[leftmargin=*]
    \item \imp{氨基甲酸酯类农药}：与有机磷类似但ChE抑制为\keyword{可逆性}，临床症状较轻。\textbf{治疗：}阿托品（\textbf{不宜使用肟类复能剂}）
    \item \imp{拟除虫菊酯类农药}：作用于神经系统钠通道。面部灼烧感、皮肤刺激、抽搐。对症支持为主
    \item \imp{百草枯（Paraquat）}：\keyword{剧毒}，特征性损害——\keyword{进行性不可逆肺纤维化}，多死于呼吸衰竭。无特效解毒剂
\end{itemize}

% ----- 3.11 复习题 -----
\section{复习题}

\begin{enumerate}[leftmargin=*, itemsep=8pt]
    \item \textbf{【名词解释】}毒物；生产性毒物；职业中毒；代谢活化
    \item \textbf{【名词解释】}中毒性肺水肿；高铁血红蛋白血症
    \item \textbf{【名词解释】}单纯窒息性气体与化学窒息性气体
    \item \textbf{【简答题】}简述毒物进入机体的主要途径及影响因素。
    \item \textbf{【简答题】}简述毒物生物转化的I相和II相反应及其意义。
    \item \textbf{【简答题】}简述铅中毒致卟啉代谢障碍的机制及实验室诊断指标。
    \item \textbf{【简答题】}简述慢性铅中毒、汞中毒、锰中毒的主要临床表现对比。
    \item \textbf{【简答题】}简述刺激性气体的毒作用特点及中毒性肺水肿的四期临床过程。
    \item \textbf{【简答题】}试比较CO和HCN的中毒机制、临床表现及解毒治疗。
    \item \textbf{【简答题】}简述苯的慢性中毒表现及其致癌性。
    \item \textbf{【简答题】}简述有机磷农药的中毒机制、临床表现及特效解毒治疗（含阿托品和复能剂的使用特点）。
    \item \textbf{【简答题】}简述百草枯中毒的特征性损害及治疗原则。
\end{enumerate}

\subsection*{参考答案要点}

\begin{enumerate}[leftmargin=*, itemsep=5pt]
    \item \textbf{毒物}：较小剂量即可引起机体病理改变甚至危及生命的化学物质。\textbf{生产性毒物}：生产过程中产生、存在于工作环境中的毒物。\textbf{职业中毒}：劳动者在生产过程中接触生产性毒物引起的中毒。\textbf{代谢活化}：无毒化学物经生物转化生成有毒中间产物的过程。
    \item \textbf{中毒性肺水肿}：吸入高浓度刺激性气体后引起肺泡及肺间质过量体液潴留的病理过程。\textbf{高铁血红蛋白血症}：Fe$^{2+}$-Hb氧化为Fe$^{3+}$-Hb失去携氧能力，以化学性青紫为特征。
    \item \textbf{单纯窒息性}：稀释空气中氧浓度导致缺氧（N$_2$、CH$_4$）。\textbf{化学窒息性}：血液窒息性（CO竞争Hb）和细胞窒息性（HCN、H$_2$S抑制细胞呼吸酶）。
    \item 三大途径：呼吸道（最主要）、消化道、皮肤。影响因素：浓度/分压、水溶性、血/气分配系数、劳动强度、肺通气量等。
    \item I相：氧化、还原、水解（官能团化，引入极性基团）。II相：结合反应（与葡萄糖醛酸、硫酸等结合）。意义：增加水溶性促进排泄，但也可代谢活化增毒。
    \item 铅抑制ALAD（→ALA↑）、亚铁螯合酶（→FEP↑、ZPP↑）→血红素合成障碍→Hb↓。诊断指标：血铅、尿铅、$\delta$-ALA↑、FEP↑、ZPP↑、点彩红细胞↑。
    \item 铅：类神经症、周围神经病（垂腕垂足）、铅绞痛、贫血、铅线。汞：易兴奋症（特有）、意向性震颤、口腔炎。锰：帕金森综合征（面具脸、慌张步态、齿轮样肌张力、静止性震颤）。
    \item 特点：病变程度取决于浓度和时间；部位与水溶性相关。四期：刺激期→潜伏期（假象期，2--6h，关键期）→肺水肿期（24h剧变，粉红色泡沫痰）→恢复期（2--15天）。
    \item CO：与Hb结合→HbCO（樱桃红色），亲和力大200--300倍。氧疗/高压氧。HCN：抑制细胞色素氧化酶Fe$^{3+}$→中断呼吸链→细胞窒息（苦杏仁味）。亚硝酸钠-硫代硫酸钠疗法。
    \item 慢性苯中毒：造血系统损害（白细胞↓→血小板↓→全血↓→再障→白血病）。IARC 1类致癌物（致白血病）。
    \item 机制：不可逆抑制AChE→ACh蓄积→胆碱能危象。M样（缩瞳、流涎、肺水肿）+N样（肌束震颤、肌麻痹）+CNS症状。解毒：阿托品（早期足量反复，抗M样）+氯磷定/解磷定（复能ChE，抗N样，对老化酶无效）。
    \item 特征性损害：进行性不可逆肺纤维化，多死于呼吸衰竭。无特效解毒剂。治疗：早期血液灌流、免疫抑制剂、对症支持。
\end{enumerate}

\newpage
% =====================================================================
%  CHAPTER 4: 粉尘与职业性肺部疾病
% =====================================================================
\newpage